\documentclass[../../main.tex]{subfiles}
\begin{document}

{\bf [解答]}

% TODO: fig: \triangle ABC

$A(0,a), B(b,0), C(-c,0) \ (a,b,c > 0)$ とおく．

\subparagraph{(イ)}

$D,E,F$は$t=0$で$A,B,C$を出発し，$t=1$で$B,C,A$に到達するから，
$D,E,F$ の時刻 $t$ での座標は
\begin{align}
  \overrightarrow{OD} &=
  \begin{pmatrix} 0 \\ a
  \end{pmatrix} + t
  \begin{pmatrix} b \\ -a
  \end{pmatrix} \\
  \overrightarrow{OE} &=
  \begin{pmatrix} b \\ 0
  \end{pmatrix} + t
  \begin{pmatrix} -c-b \\ 0
  \end{pmatrix}  \\
  \overrightarrow{OF} &=
  \begin{pmatrix} -c \\ 0
  \end{pmatrix} + t
  \begin{pmatrix} c \\ a
  \end{pmatrix} \label{eq:1}
\end{align}
となる．

従って$\triangle DEF$ の重心 $G(X,Y)$ とすると
\begin{align}
  X &= \dfrac{1}{3} [ tb + b + t(-c-b) - c + tc ] = \dfrac{1}{3} (b-c) \\
  Y &= \dfrac{1}{3} [ a - ta + ta ] = \dfrac{1}{3} a
\end{align}
となって，はじめの重心からうごかない．

\subparagraph{(ロ)}

\cref{eq:1}から
\begin{align}
  \overrightarrow{ED} &=
  \begin{pmatrix} tb \\ a(1-t)
  \end{pmatrix} -
  \begin{pmatrix} b-t(b+c) \\ 0
  \end{pmatrix} =
  \begin{pmatrix} -b+t(2b+c) \\ a(1-t)
  \end{pmatrix} \\
  \overrightarrow{EF} &=
  \begin{pmatrix} c(t-1) \\ ta
  \end{pmatrix} -
  \begin{pmatrix} b-t(b+c) \\ 0
  \end{pmatrix} =
  \begin{pmatrix} -b-c+t(b+2c) \\ ta
  \end{pmatrix}
\end{align}
だから，時刻 $t$ での $\triangle DEF$ の面積 $S(t)$ はサラスの公式から
\begin{align}
  S(t) &= \dfrac{1}{2} \left| [-b+t(2b+c)] ta - [-b-c+t(b+2c)] a(1-t) \right| \\
  &= \dfrac{1}{2} a \left| -bt + (2b+c)t^2 + (b+2c)t^2 - (2b+3c)t + b+c \right| \\
  &= \dfrac{1}{2} a \left| 3(b+c)t^2 - 3(b+c)t + b+c \right| \\
  &= \dfrac{1}{2} a(b+c) \left| 3t^2 - 3t + 1 \right| \\
  &= \dfrac{1}{2} a(b+c) \left| 3\left(t - \dfrac{1}{2}\right)^2 + \dfrac{1}{4} \right| \\
  &= \dfrac{1}{2} a(b+c) \left[ 3\left(t - \dfrac{1}{2}\right)^2 + \dfrac{1}{4} \right] \\
\end{align}
と書ける．よって$S(t)$ は $t = \dfrac{1}{2}$ の時最小値$S(1/2)=\dfrac{a(b+c)}{8}$で，この時 $\triangle ABC$ の面積 $S=\dfrac{a(b+c)}{2}$ だから
\begin{align}
  \dfrac{S(1/2)}{S} = \dfrac{1}{4}
\end{align}
である．つまり$S(1/2)$は$S$の$\dfrac{1}{4}$倍である．

\newpage

\end{document}
